\section{FreeGrasp: Visual Prompting en la sujeción robótica}

\subsection{Planteamiento del problema}

FreeGrasp resuelve el problema de la sujeción robótica de forma libre (free-form grasping): el usuario da una instrucción de sujeción ambigua en lenguaje natural (por ejemplo, «agarra el coche de juguete rojo que está abajo»), y el robot debe entender la instrucción, localizar el objeto objetivo, manejar las oclusiones y, por último, completar la sujeción.

\begin{knowledgebox}{El manejo de oclusiones es el reto central}
En escenarios reales de sujeción, el objeto objetivo suele quedar oculto por otros objetos. FreeGrasp diseñó un Pipeline completo para manejar las oclusiones: primero identifica el objetivo → determina si está oculto → retira el objeto que lo oculta → por último sujeta el objetivo. Este proceso requiere que el VLM tenga capacidad de razonamiento en varios pasos.
\end{knowledgebox}

\subsection{Diseño del Pipeline}

El flujo de procesamiento de FreeGrasp:

\begin{enumerate}
    \item \textbf{Entrada}: imagen RGBD (RGB para el razonamiento visual, Depth para la ejecución del brazo robótico) + instrucción en lenguaje natural
    \item \textbf{Detección de la imagen completa}: se usa un modelo de CV preentrenado para detectar todos los objetos de la imagen y obtener el punto central de cada uno
    \item \textbf{Visual Marking}: se marca un número en la posición de cada punto central
    \item \textbf{Razonamiento del VLM}: mediante un Prompt estructurado, se hace que el VLM razone en varios pasos a partir de la imagen con números:
    \begin{itemize}
        \item Identificar el objeto objetivo (¿qué número es el objetivo?)
        \item Determinar si está oculto (¿qué objetos ocultan al objetivo?)
        \item Planear el orden en que se retiran los objetos que lo ocultan
    \end{itemize}
    \item \textbf{Segmentación y sujeción}: se segmenta el objeto objetivo, un algoritmo tradicional estima la postura de sujeción y el brazo robótico la ejecuta
\end{enumerate}

\begin{importantbox}{El significado de Free-form}
En FreeGrasp, «Free» se refiere al grado de libertad de la entrada del usuario: se puede describir el objetivo de la sujeción con cualquier lenguaje natural, sin necesidad de un formato de instrucción estructurado. Esto se debe a la sólida capacidad de comprensión del lenguaje natural que tiene el VLM.
\end{importantbox}

\subsection{Resumen de la sección}

FreeGrasp aplica la técnica de Visual Prompting al ámbito de la sujeción robótica: mediante un Pipeline de detección → numeración → razonamiento del VLM → segmentación → ejecución, logra la sujeción de objetos a partir de instrucciones en lenguaje natural de forma libre, y la detección y el manejo de oclusiones son la innovación central.
